\section{Architecture and Framework}


Any explainable model typically terminates at the attribution layer, after finding the important 
bits. A few of them understand their causal interpretations and validate them with visual 
representations. But verifying these bits and their significance is a behavioural gap that our 
frameworks attempt to cover. To address this limitation we developed a multi-layer 
explainability framework.The neural distinguisher model, introduced by Gohr, forms the base of 
the research, on which we apply multiple layers of case studies, interpretations and validations. 

The architecture consists of the following layers: 

1.  Data Generation Layer 

2.  Neural Attribution Layer 

3.  Causal Validation Layer 

4.  Faithfulness Evaluation Layer 

5.  Interaction Analysis Layer 

6.  Structural Transport Layer 

7.  Randomized Control Layer 

8.  Structure Destruction Control Layer 

Every layer contributes to the conclusions derived from the framework using perturbation, 
ablation, transportation, destruction, randomization, causal learning and validation. The complete 
framework transforms a black box model into an interpretable system. 

\subsection{Dataset Construction and Generation Layer}

Unlike most of the research works that are based on publicly available fixed datasets, 
cryptographic data has a unique property of generating itself whenever needed, as long as the 
encryption algorithm, key generation process and input differentials are known. Here we use the 
Speck32/64 encryption algorithm to generate ciphertext pairs.  

Since the framework is based on binary classification, the algorithm produces two classes of 
ciphertext pairs. The first class contains structured samples (generated using encrypted plain 
texts), and the second class contains random samples. The Speck32/64 encrypts random plaintext 
pairs that are generated using a predefined input differential structure. The differential 
characteristics are preserved among the structured samples, while the random data loses their 
features. 


Once generated, every ciphertext pair is converted into a 64-dimensional binary feature vector. 
Since each sample consists of four 16-bit ciphertext words, the binary representation allows the 
neural network to analyze relationships directly at the bit level. This helps the ARX based 
operations to work smoothly during the encryption process as the complexity increases. 

The generated binary vectors are then classified into structured or random classes using the 
pretrained DeepSpeck neural distinguisher. The same dataset used throughout the process, in 
every layer and experiments conducted. This ensures that any observations made are due to the 
explainability layers and not the dataset itself. 

   Random Plaintext Pairs 

          ↓ 

  Introduce Differential (Structured Samples) 

          ↓ 

                Speck32/64 Encryption 

          ↓ 

                      Ciphertext Pairs 

          ↓ 

                                            Binary Conversion 

          ↓ 

         64-Dimensional Feature Vector 

          ↓ 

         DeepSpeck Neural Distinguisher 

\subsection{Attribution Layer}

The second layer of the framework aims to answer the most fundamental explainability question: 

Which input bits influence the neural network's decisions? 

To answer this question, Integrated Gradients (IG) was employed. Integrated Gradients is a 
gradient-based attribution technique that measures how much each input feature contributes to a 
model's prediction. Unlike ordinary gradients, which often suffer from instability and saturation 
effects, Integrated Gradients accumulate gradients along a continuous path between a baseline 
input and the actual sample. 

For each ciphertext sample: 

1.  A dataset of ciphertext pairs was generated using the Speck32/64 cipher. 


2.  Samples were passed through the pretrained DeepSpeck model. 

3.  Integrated Gradients was computed for each sample. 

4.  Attribution scores were calculated for all 64 input bits. 

5.  Scores were aggregated across multiple samples. 

6.  A global importance map was generated. 

This process is repeated across thousands of samples. 

The attribution scores are then averaged to obtain a global importance map. 

The resulting importance distribution reveals which regions of the ciphertext the neural network 
considers most relevant during classification. 

This layer produces the first explainability artifact of the framework: the hotspot map. 

\subsection{Causal Validation Layer}

The previous layer tells us the significant bits but doesn't give reasons that can justify their 
importance. It doesn’t prove the impact of these bits on the model’s decision.Therefore, the third 
layer introduces causal validation. 

What happens when the model is deprived of its most important input bits? 

The objective of this stage is to test whether the importance discovered by Integrated Gradients 
corresponds to genuine influence on model behavior.The procedure is straightforward. 

For each input bit: 

1.  The baseline accuracy of the neural distinguisher was recorded. 

2.  Each input bit was removed individually. 

3.  The modified samples were re-evaluated by the model. 

4.  The change in classification accuracy was measured. 

5.  Accuracy loss was recorded as the causal importance score. 


If removing a bit causes a substantial reduction in accuracy, that bit is considered causally 
important. 

If the accuracy remains unchanged, the bit is considered non-essential regardless of its attribution 
score. 

\subsection{Faithfulness Verification Layer}

Can the explanations given by the model be trusted? Do their fidelity scores 
justify their accuracy? Do the explanations generated by the model faithfully reflect 
its actual decision-making process? 

To answer the above questions, this layer was introduced. Once attribution scores and 
causal scores are obtained, the next challenge is determining whether they agree.   

This layer evaluates the faithfulness of the explanations. 

Faithfulness refers to the degree to which an explanation reflects the true behavior of the model. 

A faithful explanation should satisfy the following condition: 

High Attribution 

↓ 

        High Causal Impact 

To evaluate this relationship, the framework compares: 

●  Integrated Gradient scores 

●  Causal ablation scores 

Statistical correlation analysis is then performed between both distributions. 

Additionally, overlap analysis is conducted between the most important attribution bits and the 
most important causal bits. Quantitative approach like top-K overlap is used to measure the 
extent to which the features highlighted by the attribution method correspond to those that 
genuinely influence the model's predictions, thereby providing an indication of the faithfulness 
of the generated explanations. 

The faithfulness of explanations is decided using correlation coefficients, statistical significance 
value and top-K overlap scores. This stage is the first quantitative assessment of explanation 
quality. 


\subsection{Interaction Analysis Layer}

Cryptographic algorithms rarely depend on a single bit acting independently. Through the 
repeated mixing steps used in encryption, information from many bits becomes tightly 
intertwined as it passes through successive rounds. As a result, a neural distinguisher may 
uncover relationships that arise from multiple bits working together, rather than from any one bit 
alone. Because of this, evaluating bits in isolation can only provide a partial picture of how the 
model makes its decisions. This subsection therefore explores that issue in detail. 

Do the most important bits operate independently, or do they work together to 

influence the model's decisions? 

To investigate these hidden dependencies, an Interaction Analysis Layer was introduced after the 
causal validation stage. The primary goal of this layer is to determine whether the neural network 
relies on combinations of important bits and to identify how these bits interact with one another 
during classification. 

The following steps are performed within this layer: 

1.  The highest-ranked bits are selected pairwise 

2.  For each possible pair of important bits, both bits are simultaneously removed from the 

input data 

3.  The modified samples are then passed through the neural distinguisher 

4.  The resulting classification accuracy is recorded 

5.  This combined accuracy loss is compared against the accuracy drops observed when the 

same bits were removed individually 

6.  The interaction scores obtained from all bit pairs are aggregated into an interaction matrix 

that captures the dependency structure learned by the model 

By comparing these values, it becomes possible to measure whether the effect of removing two 
bits together is greater than, smaller than, or approximately equal to the sum of their individual 
effects.  

If removing both bits causes a substantially larger performance degradation than expected, the 
pair is considered to exhibit a cooperative or synergistic relationship. This indicates that the 
neural network relies on information shared between those bits. Conversely, if the combined 
effect is smaller than expected, the bits may contain overlapping information, suggesting 
redundancy within the learned representation. 


\subsection{Structural Stability Layer}

Do the learned cryptographic patterns remain stable as the number of encryption 

rounds increases? 

After identifying the single bit, their correlations and synergistic structures, it is important to 
understand their longevity. Can these discovered patterns survive complex ciphers?  

●  If the same regions of the ciphertext continue to influence the model across multiple 

encryption rounds, it suggests that the neural network has learned genuine cryptographic 
structures.  

●  Conversely, if the important regions change completely from one round configuration to 
another, the model may be relying on round-specific artifacts rather than fundamental 
properties of the cipher. 

To investigate this, a Structural Stability Layer was introduced within the framework. The 
purpose of this layer is to track how the model's attention changes as the number of encryption 
rounds increases and to determine whether important patterns persist throughout the encryption 
process. 

The analysis is performed in the following stages: 

1.  Multiple neural distinguishers trained on different Speck round configurations are 

selected 

2.  Every model is analyzed separately with a different level of cryptographic complexity 

3.  Every round configuration is performed independently and respect importance maps are 

generated 

4.  Each map represents how strongly the model relies on every ciphertext bit while making 

its predictions 

5.  The highest-ranked bits from each attribution map are extracted and analyzed 

6.  Correlation analysis is performed between the importance distributions obtained from 

different round configurations 

7.  The comparison is carried out between: 

●  Early and intermediate rounds 

●  Intermediate and later rounds 

●  Early and later rounds 

8.  The framework then examines how important regions evolve across rounds 


9.  We study how groups of important bits shift, expand, contract, or migrate between 

neighboring positions, leading to a phenomenon known as pattern transport 

10. The resulting correlation values are used to assess structural stability. 

●  High correlation indicates that the model continues to focus on similar regions 
across multiple rounds. This suggests that the neural distinguisher has learned 
persistent cryptographic structures rather than isolated patterns. 

●  Moderate correlation indicates partial stability, where some important regions 
remain relevant while others evolve as the cipher becomes more complex. 

●  Low correlation indicates that the model changes its attention significantly 
between round configurations. Such behavior may suggest reliance on 
round-specific features or unstable decision-making strategies. 

The Structural Stability Layer therefore provides a longitudinal view of the neural network's 
behavior. Rather than analyzing importance at a single point in time, it examines how learned 
representations evolve throughout the encryption process. This makes it possible to determine 
whether the explanations discovered by the framework represent stable cryptographic knowledge 
or temporary patterns that emerge only under specific round configurations. 

\subsection{Control Validation Layer}

The final stage of the framework consists of control experiments. We need to know that our 
explanations were based on relevant patterns and structures. Control experiments are essential 
because they test whether the discovered explanations remain valid under extreme conditions. 

Without controls, it is impossible to determine whether observed patterns represent genuine 
cryptographic structure or random artifacts. 

For this reason, two independent control experiments were performed. 

●  Randomized Control Validation 

In this experiment, the underlying structure of the input data is randomized while preserving the 
overall data format. 

The purpose is to evaluate whether attribution hotspots emerge even when meaningful 
cryptographic information is absent. 

If similar hotspot patterns appear after randomization, the explanations may be unreliable. 


Conversely, substantial differences indicate that the discovered hotspots are likely linked to 
genuine cryptographic structure. 

This experiment acts as a sanity check for the entire explainability pipeline. 

●  True Structure Destruction Control 

This experiment represents the most rigorous validation stage of the framework. 

Instead of randomizing selected components, the cryptographic structure itself is intentionally 
destroyed. 

The modified dataset retains the same dimensions and statistical format but no longer preserves 
the meaningful relationships generated by the Speck encryption process. 

The explainability pipeline is then repeated. 

If the discovered hotspots disappear or change significantly, it provides strong evidence that the 
original explanations were linked to real cipher structure. 

If the hotspots remain unchanged, the explanations may be capturing dataset artifacts rather than 
cryptographic behavior. 

This experiment provides the strongest form of validation within the framework. 

\subsection{Framework Summary}

The proposed framework progressively transforms neural cryptanalysis from a prediction task 
into an interpretable scientific investigation.The complete workflow can be summarized as: 

Data Generation 

          ↓ 

Integrated Gradients Attribution 

          ↓ 

Causal Validation 

          ↓ 

      Faithfulness Verification 

          ↓ 

         Interaction Analysis 

          ↓ 


    Structural Stability Analysis 

          ↓ 

Randomized Control Validation 

          ↓ 

Structure Destruction Validation 

↓ 

        Statistical Interpretation 

Each layer contributes a distinct form of evidence. Together, these layers provide a 
comprehensive evaluation of whether the neural distinguisher has learned genuine cryptographic 
structure and whether the generated explanations can be trusted. 

3.6. Testing methodology 

3.6.1 Experimental Setup 

The proposed explainability framework was evaluated using the publicly available DeepSpeck 
neural distinguisher for the Speck32/64 block cipher. The objective of the experimental setup 
was to provide a controlled environment for analyzing the internal decision-making process of 
the neural cryptanalysis model. 

The experimental setup consisted of the following components: 

1. Neural Model 

●  Pretrained DeepSpeck neural distinguisher 

●  Trained for ciphertext pair classification 

●  Used as the base model for all explainability experiments 

2. Programming Environment 

●  Python 

●  TensorFlow / Keras 

●  NumPy 

●  Statistical analysis libraries 

These tools were used for model execution, attribution analysis, causal interventions, and result 
evaluation. 

        

3. Dataset Generation 

●  Dataset generated directly using the Speck32/64 encryption algorithm 

●  Structured ciphertext pairs generated using differential cryptanalysis principles 

●  Random ciphertext pairs generated as the negative class 

●  Each sample converted into a 64-bit binary feature vector 

4. Experimental Consistency 

●  The same data generation procedure was maintained throughout the framework 

●  All explainability experiments were performed on a controlled input distribution 

●  This ensured fair comparison between different analysis layers 

5. Evaluation Strategy 
 The framework was not evaluated solely on prediction accuracy. Instead, multiple aspects of 
model behavior were examined: 

●  Bit-level importance 

●  Causal influence 

●  Explanation faithfulness 

●  Feature interactions 

●  Structural stability across rounds 

●  Robustness under control experiments 

By combining a benchmark neural cryptanalysis model with controlled data generation and 
multiple validation procedures, the experimental setup provides a reliable environment for 
investigating how and why the neural distinguisher makes its predictions. 

3.6.2 Attribution Evaluation Methodology 

How was feature importance identified within the neural distinguisher? 

The Attribution Layer was evaluated using the Integrated Gradients (IG) technique. The 
objective of this experiment was to identify which ciphertext bits contributed most strongly to 
the model's predictions. The results that validated this layer were: 


●  Bit-level importance scores 

●  Ranked list of influential bits 

●  Global attribution map 

This experiment provided the initial explanation of the model's behavior and served as the 
foundation for all subsequent validation stages. 

3.6.3 Causal Validation Methodology 

Are the bits identified by Integrated Gradients genuinely responsible for the model's 

decisions? 

To validate the importance scores obtained from Integrated Gradients, a causal intervention 
experiment was performed using single-bit ablation. Certain factors were responsible for the 
conclusions drawn from this layer. Like, 

●  Accuracy drop for every input bit 

●  Causal importance ranking 

●  Identification of truly influential features 

This experiment transformed attribution-based observations into experimentally verified causal 
evidence. 

3.6.4 Faithfulness Evaluation Methodology 

Do the explanations generated by the model faithfully reflect its actual 

decision-making process? 

The purpose of this experiment was to determine whether the attribution scores produced by 
Integrated Gradients aligned with the causal importance measurements obtained through 
ablation. 

Procedure: 

●  Integrated Gradients scores were collected. 

●  Causal importance scores were collected. 

●  Correlation analysis was performed between the two rankings. 


●  Top-k overlap analysis was conducted. 

●  Agreement scores were calculated. 

Evaluation Measures: 

●  Pearson Correlation 

●  Top-10 Overlap 

●  Top-20 Overlap 

●  Agreement Score 

Higher agreement indicates that the generated explanations accurately reflect the model's 
behavior. 

3.6.5 Interaction Analysis Methodology 

This experiment investigated whether the neural distinguisher relied on isolated features or 
groups of interacting features. Certain conclusions were drawn using the generated measures 
below, 

●  Pairwise interaction matrix 

●  Synergy scores 

●  Redundancy scores 

The experiment revealed hidden dependencies among important features learned by the model. 

3.6.6 Structural Stability Evaluation Methodology 

This experiment examined whether the patterns identified by the model persisted across different 
levels of encryption complexity. Comparisons performed were evaluated using, 

●  Round-to-round correlation scores 

●  Pattern transport measurements 

●  Stability indicators 

This experiment assessed whether the model had learned persistent cryptographic structures or 
round-specific patterns. 

 
 

3.6.7 Control Validation Methodology 

Control experiments were conducted using randomized data. Below given are the steps it 
followed to conduct this experiment. 

●  Attribution maps were generated using genuine ciphertext samples. 

●  A second dataset with randomized structure was created. 

●  The same explainability procedures were applied to both datasets. 

●  Importance distributions were compared. 

●  Entropy and similarity measures were calculated. 

Evaluation Measures: 

●  Entropy 

●  Correlation 

●  Importance Distribution Similarity 

This experiment served as a sanity check and verified whether the framework was identifying 
meaningful cryptographic patterns rather than random noise. 

4. RESULTS AND OUTCOMES 
